\section{Arrays}
\label{sec:arrays}

Arrays are the language's one growable sequence. They are
\textbf{reference values}: a binding names an array \emph{object}, and
mutating operators write into that object so every binding sees the
change. Arrays are \textbf{homogeneous}: one element type, enforced
everywhere.

\subsection{Structure: length and capacity}
\label{sec:lencap}

An array object carries two independent quantities:

\begin{center}
\begin{tabular}{lll}
\toprule
 & meaning & invariant \\
\midrule
\ic{len} & number of live elements & $\texttt{len} \ge 0$ \\
\ic{cap} & elements the storage can hold & $\texttt{len} \le \texttt{cap}$ \\
\bottomrule
\end{tabular}
\end{center}

\noindent Indexing is trapped: reading or writing \ic{a[i]} with
\ic{i < 0} or \ic{i >= len} is a runtime error (``index out of
bounds''). \ic{cap} is never observable through operators --- only
\ic{len} is (\ic{len(a)}).

\paragraph{Growth policy (amortized doubling).}
\begin{align*}
\texttt{cap}_0 &= 0 \\
\texttt{cap}_{k+1} &= \max(8,\; 2\cdot\texttt{cap}_k)
   \qquad \text{when } \texttt{len} \text{ reaches } \texttt{cap}
\end{align*}

\noindent That is: the array grows \emph{only} when the capacity is
reached; the first growth yields capacity 8 and every later growth
doubles it ($8, 16, 32, \ldots$). Appending $n$ elements one at a time
therefore costs $O(n)$ copies total, $O(1)$ amortized per element. The
capacity never shrinks.

\subsection{Operators}
\label{sec:array-ops}

Let \ty{a}, \ty{b} be arrays of \ty{T}, and \ty{x} a value of type
\ty{T}. The complete operator semantics:

% The four columns are wide; a tighter \tabcolsep keeps the table inside
% the text block.
\begingroup
\setlength{\tabcolsep}{3.4pt}
\begin{center}
\begin{tabular}{llll}
\toprule
expression & result & effect on operands & value types \\
\midrule
\ic{a + b} & new array $= a \mathbin{+\!\!+} b$ & none &
   \ty{T} $=$ \ty{T} \\
\ic{a + x} & new array $= a \mathbin{+\!\!+} [\,x\,]$ & none &
   type of \ty{x} is \ty{T} \\
\ic{x + a} & new array $= [\,x\,] \mathbin{+\!\!+} a$ & none &
   type of \ty{x} is \ty{T} \\
\ic{rows + row} & row appended whole & none & \ic{[[T]]}, \ic{[T]} \\
\ic{row + rows} & row prepended whole & none & \ic{[[T]]}, \ic{[T]} \\
\ic{x :: a} & new array $= [\,x\,] \mathbin{+\!\!+} a$ & none &
   \ty{x} is ONE element \\
\ic{a += b} & \ty{a} & appends every element of \ty{b}, in place &
   \ty{T} $=$ \ty{T} \\
\ic{a += x} & \ty{a} & pushes \ty{x}, in place &
   type of \ty{x} is \ty{T} \\
\ic{a *= x} & \ty{a} & prepends \ty{x}, in place &
   type of \ty{x} is \ty{T} \\
\ic{a = b}  & \ty{b} & aliases: both names name one array & \\
\bottomrule
\end{tabular}
\end{center}
\endgroup

\begin{sloppypar}
\noindent The two \ic{rows} rows are the array-as-an-element case, and
they say what \ic{[[T]]} plus a \ic{[T]} does: the array goes in
\emph{whole}, as one element, because a \ic{[T]} is what that array
holds. \ic{rows + rows} (both \ic{[[T]]}) concatenates as usual, and so
does \ic{rows += row}, which appends the row in place.
\end{sloppypar}

\noindent In prose, matching the design intent exactly:

\begin{itemize}[nosep]
\item \ic{arr1 + arr2} --- make a new array which is \ic{arr1}
  concatenated with \ic{arr2}; the operands are untouched.
\item \ic{arr1 += arr2} --- append \ic{arr2} to the end of \ic{arr1},
  \emph{in place}.
\item \ic{arr1 = arr2} --- both point to the same array in memory.
\item \ic{arr1 + elem} --- a new array which is \ic{arr1} concatenated
  with \ic{[elem]}; only allowed if \ic{elem} has, or converts to, the
  element type (Section~\ref{sec:conversions}).
\item \ic{elem + arr1} --- a new array which is \ic{[elem]} concatenated
  with \ic{arr1}.
\item \ic{+=} means in-place push (of an element) or append (of an
  array); \ic{+} always means a \emph{new} array; \ic{*=} means prepend
  the item in place. A push and a prepend both take amortized constant
  time: an array keeps spare room at both ends and doubles it when it
  runs out, so building an array from either end is linear.
\end{itemize}

\paragraph{The one-thing rule.} Whether \ic{+=} pushes or appends is
decided by the \emph{shape of the value against the target's element
type}: an array whose own elements are the target's element type has
those elements appended; anything else --- a single value, or an array of
some other element type --- is pushed as one element. So a \ic{[long]}
onto a \ic{[[long]]} is pushed whole (it is one \ic{[long]} element), a
\ic{[[long]]} onto a \ic{[[long]]} appends its rows, and a \ic{[str]} is
neither, so it is an error:

\begin{lstlisting}
let rows = [[1], [2]];
rows += [3];        // pushes the one row [3] -- three rows now
rows += [[4], [5]]; // appends two rows -- five rows now
rows += 6;          // error: '+=' element: expected [long], got long
rows += ["s"];      // error: '+=' appended elements: expected [long], got str
\end{lstlisting}

\paragraph{Homogeneity.} Every mixing form requires matching element
types; the checker rejects mismatches --- including the append case, where
a non-matching array would splice wrong-typed elements:

\begin{lstlisting}
let xs = [1, 2];
let fs: [float] = [0.5];
xs += fs;           // error: '+=' appended elements: expected long, got float
\end{lstlisting}

\paragraph{\ic{a += a}.} The right side is captured (its length snapshotted)
before appending, so \ic{a += a} doubles \ic{a} with a copy of its
original elements.

\subsection{Aliasing and visibility}
\label{sec:array-aliasing}

Assigning arrays aliases; mutating operators are visible through every
alias:

\begin{lstlisting}
let m1 = [1, 2];
let m2 = m1;      // ONE array, two names
m1 += 3;          // m2 sees it: len(m2) == 3
\end{lstlisting}

\noindent By contrast \ic{+} never mutates: \ic{let ys = xs + [9];}
leaves \ic{xs} at its old length.

\subsection{Fixed arrays}

\ic{[T; N]} arrays are allocated with exactly \ic{N} slots; their length
is part of the type, so they reject growth: \ic{+=} and \ic{*=} on a
fixed array is a compile error (``cannot grow fixed array \ic{[long; 3]}
with \ic{'+='}''). Indexing is bounds-trapped like heap arrays. A fixed
array may be sliced (Section~\ref{sec:slices}) and slice-assigned
(below); the result of a slice is a heap \ic{[T]}.

Since the length is in the type, an annotated fixed array checks what it
is given: a heap \ic{[T]} never converts to a \ic{[T; N]}, and a literal
is checked against the size it claims --- \ic{let v: [long; 5] = [1, 2];}
is an error (\ic{expected [long; 5], got [long]}), while the matching
\ic{[long; 2]} is fine. Build one with \ic{[v; n]}, which is a fixed array
of \ty{n} copies of \ic{v} (Section~\ref{sec:expr-grammar}).

\subsection{Slices}
\label{sec:slices}

\ic{a[lo:hi]} evaluates to a \textbf{new} heap \ic{[T]} holding the
elements in $[\texttt{lo}, \texttt{hi})$. Missing bounds default to
\ic{0} and \ic{len}; bounds are clamped (no trap). All four spellings
are legal: \ic{a[:]}, \ic{a[lo:]}, \ic{a[:hi]}, \ic{a[lo:hi]}. Strings
slice the same way: \ic{s[lo:hi]} is a new \ic{str} of the bytes in that
range, clamped alike. There is no \ic{substr}.

\begin{sloppypar}
\paragraph{Slice assignment.} \ic{a[lo:hi] = b} writes \ic{b} into \ic{a}
in place, over the range \ic{[lo, hi)} --- the same clamped range a read
would take, on either kind of array. It is a statement, so as the last
one of a block it is that block's value (Section~\ref{sec:assignment}).
The target's subexpressions, the bounds and \ic{b} are each evaluated
once, left to right (Section~\ref{sec:exactly-once}). The right side
must be an array whose element type is the target's, a \ic{str} cannot
be slice-assigned, and the length is checked when the program runs: a
\ic{b} of another length traps with
\ic{runtime error: slice assignment length mismatch}, and a range
outside the array with \ic{runtime error: slice assignment out of
bounds}. The static failures name themselves:

\begingroup
\setlength{\tabcolsep}{4pt}
\begin{center}
\begin{tabular}{ll}
\toprule
mistake & message \\
\midrule
a \ic{[float]} into a \ic{[long]} &
  \ic{slice assignment: cannot store [float] into [long]} \\
not an array &
  \ic{slice assignment needs an array on the right, got long} \\
a \ic{str} target &
  \ic{cannot assign to a slice of str (strings are immutable)} \\
\bottomrule
\end{tabular}
\end{center}
\endgroup
\end{sloppypar}

\begin{lstlisting}
let a = [1, 2, 3, 4, 5];
a[1:3] = [9, 8];            // a is [1, 9, 8, 4, 5]
a[1:3] = [7];               // traps: slice assignment length mismatch
\end{lstlisting}

\subsection{Byte strings}

A \ic{[byte]} array \emph{is} a binary string: \ic{b"..."} literals
produce them, \ic{+} concatenates, slices cut, and printing one shows
the Python-style repr --- printable ASCII verbatim, quote and backslash
escaped, \ic{\bs n} \ic{\bs r} \ic{\bs t} as escapes, and every other
byte as \ic{\bs xHH}:

\begin{lstlisting}
print(b"\bs x7fELF");     // b'\bs x7fELF'
let hdr = b"AB" + b"\bs x00";
print(hdr);              // b'AB\bs x00'
print(hdr[0:2]);         // b'AB'
\end{lstlisting}

\noindent \ic{\bs{std/struct.y}} builds on this with pack/unpack
helpers for file headers and binary layouts (\ic{u32be}, \ic{get\_i16le},
\ic{pack}, \ic{unpack}, \ldots).

\subsection{Operators are the interface}

There are no array mutator functions: \ic{+=} and \ic{*=} \emph{are} push
and prepend, and there is no \ic{push}, \ic{prepend}, \ic{map\_get} or
\ic{map\_set} to call instead (Section~\ref{sec:builtin-pointers}).
\ic{len(a)} is the only builtin whose argument is an
array \emph{as such}; \ic{join(xs, sep)} also takes an array but of any
element type (each element is stringified), and \ic{bytes}/\ic{from\_bytes}
and \ic{fs.write\_bytes} take a \ic{[byte]}.
